% ==============================================================================
% Capítulo 5: Projeto do sistema
% ==============================================================================
\chapter{Projeto do sistema}
\label{cap:projeto}

\begin{instrucao}
\orientacao{O que vai no Projeto do Sistema}
Este capítulo mostra como o sistema foi organizado para atender aos
requisitos. As competências do TCC pedem a definição de uma arquitetura
adequada, modelos de estrutura e de comportamento e a análise do projeto sob a
perspectiva de atributos de qualidade. Inclua:
\begin{itemize}
  \item a \textbf{arquitetura}: componentes, camadas, serviços externos e como
        se comunicam;
  \item o \textbf{modelo estrutural}: diagrama de classes ou modelo de dados;
  \item o \textbf{modelo comportamental}: diagramas de sequência, de atividades
        ou de estados para os fluxos mais importantes;
  \item as \textbf{decisões de projeto}: alternativas consideradas e por que a
        escolhida foi adotada, relacionando-as aos requisitos.
\end{itemize}
Não inclua diagramas sem explicação. Cada diagrama deve ser apresentado no
texto, que destaca o que ele mostra de relevante. Diagramas grandes demais
para uma página devem ser divididos ou ir para um apêndice.
\end{instrucao}

Este capítulo apresenta o projeto do Monitora. A
\cref{sec:projeto:arquitetura} descreve a arquitetura, a
\cref{sec:projeto:modelo} apresenta o modelo de classes e a
\cref{sec:projeto:decisoes} discute as principais decisões de projeto.

% ------------------------------------------------------------------------------
\section{Arquitetura}
\label{sec:projeto:arquitetura}

O sistema segue uma arquitetura em camadas, apresentada na
\cref{fig:arquitetura}. O servidor web Nginx recebe as requisições e as
encaminha à aplicação Django, executada pelo servidor de aplicação Gunicorn. A
aplicação está dividida em três módulos, que correspondem aos perfis de
usuário e aos grupos de requisitos. O módulo \texttt{contas} trata da
autenticação e dos perfis, o módulo \texttt{agenda} trata dos horários e das
reservas (RF01 a RF05) e o módulo \texttt{relatorios} gera os relatórios dos
professores (RF07). Os lembretes (RF08) são enviados por uma tarefa agendada,
executada fora do ciclo das requisições para não atrasar as respostas.

\begin{figure}[htbp]
  \centering
  \caption{Arquitetura do Monitora}
  \label{fig:arquitetura}
  \begin{tikzpicture}[
      font=\footnotesize,
      comp/.style={draw, rounded corners=2pt, minimum height=8mm,
                   minimum width=2.6cm, align=center, fill=white},
      seta/.style={-{Stealth}, thick},
    ]
    \node[comp] (nav) {Navegador};
    \node[comp, below=7mm of nav] (nginx) {Nginx};
    \node[comp, below=7mm of nginx] (gunicorn) {Gunicorn};
    \node[comp, below=10mm of gunicorn, xshift=-3.2cm] (contas) {\texttt{contas}};
    \node[comp, below=10mm of gunicorn] (agenda) {\texttt{agenda}};
    \node[comp, below=10mm of gunicorn, xshift=3.2cm] (relat) {\texttt{relatorios}};
    \node[comp, below=12mm of agenda, fill=gray!10] (bd) {PostgreSQL};
    \node[comp, right=12mm of relat] (tarefa) {Tarefa agendada\\(lembretes)};
    \begin{scope}[on background layer]
      \node[draw, dashed, rounded corners=3pt, fill=gray!5, inner sep=3mm,
            fit=(contas)(agenda)(relat),
            label={[anchor=north west, font=\scriptsize]north west:Aplicação Django}] (django) {};
      \node[draw, dotted, rounded corners=3pt, inner sep=4mm,
            fit=(nginx)(django)(bd)(tarefa),
            label={[anchor=south east, font=\scriptsize]south east:Contêineres Docker}] {};
    \end{scope}
    \draw[seta] (nav) -- node[right, font=\scriptsize] {HTTPS} (nginx);
    \draw[seta] (nginx) -- (gunicorn);
    \draw[seta] (gunicorn) -- (django.north);
    \draw[seta] (contas) |- (bd);
    \draw[seta] (agenda) -- (bd);
    \draw[seta] (relat) |- (bd);
    \draw[seta] (tarefa.south) |- ([yshift=-2.5mm]bd.east);
    \draw[seta] (tarefa.east) -- ++(9mm,0) node[right, align=left, font=\scriptsize] {servidor\\de e-mail};
  \end{tikzpicture}
  \fonteautor
\end{figure}

% ------------------------------------------------------------------------------
\section{Modelo de classes}
\label{sec:projeto:modelo}

A \cref{fig:classes} apresenta o diagrama de classes de domínio, em notação
\ac{uml}, com os atributos principais. Uma \texttt{Disciplina} tem vários
monitores e cada \texttt{Monitor} publica \texttt{Horarios} semanais. Cada
ocorrência de um horário em uma data específica pode receber no máximo uma
\texttt{Reserva}, feita por um \texttt{Estudante}. Quando o atendimento
acontece, o monitor registra o assunto e a presença na própria reserva, o que
evita uma entidade adicional e simplifica os relatórios. As classes
\texttt{Monitor}, \texttt{Estudante} e \texttt{Professor} são perfis de um
mesmo usuário do sistema de autenticação do Django, o que permite que um
monitor também reserve atendimentos em outra disciplina.

\begin{figure}[htbp]
  \centering
  \caption{Diagrama de classes de domínio do Monitora}
  \label{fig:classes}
  \begin{tikzpicture}[
      font=\scriptsize,
      classe/.style={draw, rectangle split, rectangle split parts=2,
                     rectangle split part fill={gray!15,white},
                     text width=3.3cm, align=left, inner sep=3pt},
      assoc/.style={thick},
    ]
    \node[classe] (disc) {\centering\textbf{Disciplina}
      \nodepart{second} codigo: str\\ nome: str\\ semestre: str};
    \node[classe, right=2.2cm of disc] (mon) {\centering\textbf{Monitor}
      \nodepart{second} usuario: Usuario\\ disciplina: Disciplina};
    \node[classe, below=1.6cm of mon] (hor) {\centering\textbf{Horario}
      \nodepart{second} dia\_semana: int\\ inicio: time\\ fim: time};
    \node[classe, below=1.6cm of hor] (res) {\centering\textbf{Reserva}
      \nodepart{second} data: date\\ assunto: str\\ presente: bool};
    \node[classe, left=2.2cm of res] (est) {\centering\textbf{Estudante}
      \nodepart{second} usuario: Usuario\\ matricula: str};
    \node[classe, below=1.6cm of disc] (prof) {\centering\textbf{Professor}
      \nodepart{second} usuario: Usuario\\ siape: str};
    \draw[assoc] (disc) -- node[above, pos=0.15] {1} node[above, pos=0.85] {1..*} (mon);
    \draw[assoc] (prof) -- node[left, pos=0.2] {1..*} node[left, pos=0.8] {1..*} (disc);
    \draw[assoc] (mon) -- node[right, pos=0.2] {1} node[right, pos=0.8] {0..*} (hor);
    \draw[assoc] (hor) -- node[right, pos=0.2] {1} node[right, pos=0.8] {0..*} (res);
    \draw[assoc] (est) -- node[above, pos=0.15] {1} node[above, pos=0.85] {0..*} (res);
  \end{tikzpicture}
  \fonteautor
\end{figure}

% ------------------------------------------------------------------------------
\section{Decisões de projeto e atributos de qualidade}
\label{sec:projeto:decisoes}

A decisão mais importante do projeto foi como garantir que um horário não seja
reservado por dois estudantes ao mesmo tempo (RF03). A primeira alternativa
considerada foi verificar a disponibilidade na visão antes de gravar a
reserva. Essa verificação, porém, não impede que duas requisições simultâneas
passem pela checagem antes de qualquer uma gravar. Por isso, adotou-se uma
restrição de unicidade no banco de dados sobre o par (horário, data), de modo
que a segunda gravação falhe e a aplicação informe ao estudante que o horário
acabou de ser ocupado. O \cref{cod:conflito}, no \cref{cap:implementacao},
mostra a implementação dessa regra.

Do ponto de vista dos atributos de qualidade, a separação em módulos favorece
a manutenibilidade (RNF05), porque cada módulo pode ser testado de forma
isolada. A execução dos lembretes em tarefa separada contribui para o
desempenho (RNF03). O uso do sistema de autenticação e de permissões do
próprio Django reduz o risco de falhas de segurança em código escrito do zero
(RNF04).

\begin{instrucao}
\orientacao{Decisões de projeto}
Bancas valorizam o raciocínio por trás das escolhas. Para cada decisão
relevante, apresente o problema, as alternativas consideradas, a escolhida e o
motivo, relacionando-a aos requisitos e aos atributos de qualidade. Diagramas
de sequência ajudam a explicar fluxos com várias etapas, como o de reserva
descrito acima. Se o diagrama ficar grande, mostre apenas o fluxo principal.
\end{instrucao}
